% Einheit 1 von 5 – Körper erkennen, Netze, Schrägbilder (bank/koerper/e1.jsonl, zusammenbau v0.1)
%% TODO Zweigzeile Teil 1: was hier gelernt wird (Satz in Schülersprache aus dem Einheitstitel) – Einheit 1
\einheitenkopf[e1]{Einheit 1 von 5 · Körper erkennen, Netze, Schrägbilder}
\zweigzeile{ab Kl. 5, je nach Buch bis Kl. 6 · P10 oft}
\setcounter{aufgabe}{8}

%% TODO Titel: Ich-kann-Satz fehlt in der Bank (Platzhalter: Kettenname) – Nr. 9
%% TODO Anweisung: ein Satz über der ersten Teilaufgabe fehlt in der Bank – Nr. 9
\begin{aufgabe}{Körper und Netze}
%% TODO \rechenplatz{n}: Schrittzahl des Grundfalls fehlt in der Bank (nur bei fester Schreibform, 2.3 b)
\begin{teile}
\teil Ein Zelt hat diese Form. Kreuze an, welcher Körper es ist \hfill (P10 2017 OS). \\ \kreuz{Pyramide} \\ \kreuz{Prisma} \\ \kreuz{Quader}

\prismadreieck{4}{2}{2.5}{}{}{}
\teil Ein Schuhkarton hat diese Form. Wie heißt der Körper?

\quader{4}{2}{2.5}{}{}{}
\leerfeld
\teil Ein Karton hat diese Form. Wie viele Ecken, Kanten und Flächen hat er?

\quader{4}{2}{2.5}{}{}{}
Ecken \leerfeld , Kanten \leerfeld , Flächen \leerfeld
\teil Lässt sich dieses Netz zu einem Würfel falten? \janein

$\begin{array}{cccc}  & \square &  &  \\ \square & \square & \square & \square \\  &  & \square &  \end{array}$
\teil Aus dem Netz wird ein Würfel gefaltet. Welche Fläche liegt der Fläche B gegenüber?

$\begin{array}{cccc}  & A &  &  \\ B & C & D & E \\  & F &  &  \end{array}$
\leerfeld
\teil Zeichne das Netz einer Streichholzschachtel: Quader mit $4$ cm Länge, $2$ cm Breite und $1$ cm Höhe.

\rechenplatz{5}
\teil Zeichne auf Kästchenpapier das Schrägbild eines Quaders: Länge $6$ cm, Tiefe $4$ cm, Höhe $3$ cm.

\rechenplatz{5}
\teil Ein Glasdach hat die Form einer Pyramide mit quadratischer Grundfläche, $2{,}4$ m hoch, Grundkante $3{,}6$ m. Zeichne die Höhe in das Schrägbild ein und beschrifte die Skizze mit beiden Maßen \hfill (P10 2015 OS).

\pyramide{3.5}{3.5}{3}{}{}{}
\teil Aus dem Netz wird eine Geschenkschachtel in Würfelform gefaltet. Die Fläche D ist der Deckel. Gib an, welche Fläche der Boden ist \hfill (P10 2016 OS).

$\begin{array}{cccc}  &  & A &  \\ D & N & M &  \\  &  & B & C \end{array}$
\leerfeld
\end{teile}
\end{aufgabe}

%% TODO Titel: Ich-kann-Satz fehlt in der Bank (Platzhalter: Kettenname) – Nr. 10
%% TODO Anweisung: ein Satz über der ersten Teilaufgabe fehlt in der Bank – Nr. 10
\begin{aufgabe}{Grund- und Deckfläche, Seitenflächen benennen (auch beim liegenden Prisma)}
\begin{teile}
\teil Das Prisma liegt auf einer Rechteckfläche. Welche Form haben Grund- und Deckfläche, welche die Seitenflächen?

\prismadreieck{4}{2.5}{2}{}{}{}
\end{teile}
\end{aufgabe}

%% TODO Titel: Ich-kann-Satz fehlt in der Bank (Platzhalter: Kettenname) – Nr. 11
%% TODO Anweisung: ein Satz über der ersten Teilaufgabe fehlt in der Bank – Nr. 11
\begin{aufgabe}{Schrägbild lesen: verdeckte Kanten, Maße entnehmen}
\begin{teile}
\teil Im Schrägbild eines Quaders ist die Vorderkante $5$ cm und die schräge Kante nach hinten $3$ cm lang. Wie tief ist der Körper wirklich, und wie viele Kanten sind gestrichelt? Tiefe \leerfeld[cm] , gestrichelt \leerfeld Kanten
\end{teile}
\end{aufgabe}

%% TODO Titel: Ich-kann-Satz fehlt in der Bank (Platzhalter: Kettenname) – Nr. 12
%% TODO Anweisung: ein Satz über der ersten Teilaufgabe fehlt in der Bank – Nr. 12
\begin{aufgabe}{Körper in ein Schrägbild einzeichnen (Kegel im Quader)}
\begin{teile}
\teil Ein Kegel aus Beton (Radius $25$ cm, Höhe $70$ cm) wird in einem quaderförmigen Karton verschickt, der so klein wie möglich ist. Skizziere den Kegel in das Schrägbild des Kartons \hfill (P10 2024 OS).

\quader{2.5}{2.5}{3.5}{}{}{}
\end{teile}
\end{aufgabe}

%% TODO Titel: Ich-kann-Satz fehlt in der Bank (Platzhalter: Kettenname) – Nr. 13
%% TODO Anweisung: ein Satz über der ersten Teilaufgabe fehlt in der Bank – Nr. 13
\begin{aufgabe}{Körper und Netze – Fehler finden, Begründen, Anwendung, Darstellungswechsel}
\begin{teile}
\teil Mia sagt: „Das ist eine Pyramide, weil Dreiecke daran sind.“ Wo ist der Fehler?

\prismadreieck{4}{2.5}{2}{}{}{}
\teil Warum lässt sich dieses Netz nicht zu einem Würfel falten?

$\begin{array}{ccc}  & \square &  \\ \square & \square & \square \\  & \square &  \end{array}$
\teil Für ein Spiel wird ein Würfel gebastelt. Auf einem Spielwürfel ergeben die Augenzahlen gegenüberliegender Flächen zusammen $7$. Welche Zahlen gehören in die Felder X, Y und Z?

$\begin{array}{cccc}  & 1 &  &  \\ 2 & 3 & X & Y \\  & Z &  &  \end{array}$
X = \leerfeld , Y = \leerfeld , Z = \leerfeld
\teil Aus diesem Netz wird ein Körper gefaltet. Kreuze an, welcher Körper entsteht \hfill (P10 2018 OS). \\ \kreuz{Prisma} \\ \kreuz{Pyramide} \\ \kreuz{Quader}

\netzpyramide{2.5}{2}{}{}
\end{teile}
\end{aufgabe}
